\subsection{数据证据分层与四问之间的数值接口}
同一附件中的数字承担不同角色，必须先区分“构造评分”“拟合函数”“检验预测”和“设定情景”。例如，A1 的字段可用于定义质量分数，但没有直接给出按该分数筛选后的训练损失；B1 有参数量、训练量和损失，却没有与 A1 同标度的质量标签。将两组数据接入同一个模型，需要显式保留连接假设。本文据此将主要证据分为表~\ref{tab:evidence-levels} 所列的四类。

\begin{table}[!htbp]
\centering\songti\zihao{-4}
\caption{主要附件的证据角色与使用边界}
\label{tab:evidence-levels}
\begin{tabular}{@{}>{\raggedright\arraybackslash}p{2.5cm}>{\raggedright\arraybackslash}p{4.2cm}>{\raggedright\arraybackslash}p{4.2cm}>{\raggedright\arraybackslash}p{3.5cm}@{}}
\toprule
数据类别 & 本文用途 & 可支持的判断 & 不直接支持的判断 \\
\midrule
A1、A2、A3 & 固定参照评分与同名域对照 & 相对评分、冲突发生率 & 评分提高必然降低训练损失 \\
B1、B2、B4、B5 & 经典律拟合与分层验证 & 同族拟合及跨来源误差 & 任意架构上的通用参数 \\
B6--B8 & 质量项的条件预测比较 & 半合成设置内的误差差异 & A1 质量指数的因果估计 \\
估算外推表与趋势情景 & 一致性检查、条件配置和外推 & 假设成立时的数值后果 & 真实未来性能或无条件预测 \\
\bottomrule
\end{tabular}
\end{table}

数值接口按“输入标度一致、单位一致、用途明确”的顺序建立。问题一输出样本评分、域均值和域级中位数；问题三只读取其中的 $Q_0$ 作为质量成本起算点。问题二输出的 B1 经典参数用于资源配置，而半合成经验拟合用于检验质量字段的预测信息，两套参数不拼接。问题四的损失--能力桥接与月度前沿预测也分别计算，前者检验损失和能力之间是否存在可用映射，后者直接描述榜单能力随提交月份的变化。

为使各问结论可追溯，本文对数值结果区分三种来源：从附件直接统计得到的量，例如样本数与域均值；从指定模型估计得到的量，例如经典律系数；从给定参数和情景推导得到的量，例如预算最优配置与两年前沿分位数。第一类也会受抽样和字段定义影响，第二类还会受函数形式影响，第三类则进一步依赖外推假设。后文每次比较均保持数据标度与检验角色一致，不把模型内部的一致性检查当作外部预测验证。

这种分层也决定了问题之间的解释方向。得到较高的相对质量分数，表示在当前参照与权重下更靠近所选指标的有利端；它不能直接转换成提升多少基准分数。预算模型给出的是条件损失最小的配置，仍需在实际训练中确认其损失；即使损失得到验证，能力收益也应在统一基准上再检验。由此形成“数据评分--损失模型--资源配置--能力验证”的完整研究路线，同时明确本题附件已经覆盖与尚未覆盖的环节。
